% =========================================================
%  config/styles.tex
%  Warna, tipe kolom, format judul, header/footer, keterangan tabel
% =========================================================

% --- Warna institusi -----------------------------------
\definecolor{ubPrimary}{HTML}{1F3864}
\definecolor{ubAccent}{HTML}{2E75B6}
\definecolor{ubHead}{HTML}{D9E2F3}
\definecolor{ubAlt}{HTML}{F2F6FC}
\definecolor{ubLine}{HTML}{BFBFBF}

% --- Tipe kolom ----------------------------------------
\newcolumntype{L}[1]{>{\RaggedRight\arraybackslash}p{#1}}
\newcolumntype{C}[1]{>{\Centering\arraybackslash}p{#1}}
\newcolumntype{J}[1]{>{\justifying\arraybackslash}p{#1}}
\newcolumntype{Y}{>{\RaggedRight\arraybackslash}X}

% --- Jarak antar paragraf ------------------------------
\setlength{\parindent}{0pt}
\setlength{\parskip}{0.45em}
\renewcommand{\arraystretch}{1.25}

% --- Toleransi pemenggalan baris -----------------------
% Memberi keluaran satu kali lagi dengan peregangan tambahan sehingga
% kotak teks yang sedikit melebihi lebar kolom dapat dihindari.
\setlength{\emergencystretch}{2em}
\tolerance=1500

% --- Format judul section ------------------------------
\titleformat{\section}
  {\normalfont\large\bfseries\color{ubPrimary}}
  {\thesection.}{0.5em}{}
\titlespacing*{\section}{0pt}{1.5em}{0.6em}

\titleformat{\subsection}
  {\normalfont\normalsize\bfseries\color{ubAccent}}
  {\thesubsection}{0.5em}{}
\titlespacing*{\subsection}{0pt}{1.1em}{0.4em}

\titleformat{\subsubsection}
  {\normalfont\normalsize\bfseries}
  {\thesubsubsection}{0.5em}{}
\titlespacing*{\subsubsection}{0pt}{0.9em}{0.3em}

% --- Header & footer -----------------------------------
\pagestyle{fancy}
\fancyhf{}
\fancyhead[L]{\footnotesize\color{ubPrimary}Universitas Bale Bandung}
\fancyhead[R]{\footnotesize\color{ubPrimary}RPS Pengantar Pemrograman Game}
\fancyfoot[C]{\footnotesize\thepage}
\renewcommand{\headrulewidth}{0.4pt}
\renewcommand{\footrulewidth}{0pt}

% --- Keterangan tabel & gambar -------------------------
\captionsetup{
  font=small,
  labelfont={bf,color=ubPrimary},
  skip=4pt
}

% --- Daftar (list) -------------------------------------
\setlist[itemize]{topsep=2pt,itemsep=1pt,parsep=0pt,leftmargin=1.2em}
\setlist[enumerate]{topsep=2pt,itemsep=1pt,parsep=0pt,leftmargin=1.4em}

% --- Tabel ringkas (keterangan di bawah tabel) ---------
\newcommand{\tabsource}[1]{%
  \par\vspace{2pt}{\footnotesize\itshape #1\par}%
}

% --- Garis tabel (booktabs) ----------------------------
% Header tabel ditandai dengan garis tebal di atas dan tipis di bawah,
% sebagai pengganti pewarnaan latar baris (colortbl tidak tersedia).
\renewcommand{\toprule}{\specialrule{1.3pt}{0pt}{0pt}}
\renewcommand{\midrule}{\specialrule{0.5pt}{0pt}{0pt}}

% --- Tampilan tautan -----------------------------------
\hypersetup{
  colorlinks = true,
  linkcolor  = ubPrimary,
  urlcolor   = ubAccent,
  citecolor  = ubAccent,
  filecolor  = ubAccent
}
